\clearpage
\phantomsection
\addcontentsline{toc}{chapter}{Something is afoot}
\includepdf{Illustrations/blackboard.jpg}


% \section*{Conference Dinner}

\lettrine{A}{conference} dinner, measured glasses clinking, measured hushed conversations. Ernest Olber sits beside Demi-Moorish, unaware he is scanning the crowd beyond her.

\smallskip

\textit{Demi}  
\textit{(teasing, nudging him)}  
\textit{"You’re brooding again."}

\smallskip

\textit{Ernest}  
\textit{(with a tight smile)}  
\textit{"Am I?"}

\smallskip

\textit{Demi}  
\textit{"You only do that when you see ghosts."}

\smallskip

\noindent
His grip tightens on his glass. Across the room, Eleanor and Willard are deep in conversation.

\textit{Demi}  
\textit{(following his gaze)}  
\textit{"You always had a soft spot for Eleanor."}

\begin{figure}[h]
    \centering
    \includegraphics[width=0.8\textwidth]{Illustrations/vign-DemiErnest.png}
    \caption{Ernest being so.}
    \label{fig:ernest_demi}
\end{figure}

\textit{Ernest}  
\textit{(chuckling dryly)}  
\textit{"I had a weakness for wrecking things."}

\smallskip

\textit{Demi}  
\textit{"What are you talking about?"}

\smallskip

He doesn’t answer right away. Instead, he downs his drink and murmurs, \textit{"Nothing."} But his eyes flick toward Sylvia, sitting at the far end of the room, wheelchair-bound, speaking animatedly with a colleague. Demi’s gaze follows his, and for the first time, a lurch.

% ============ COMMENTED OUT (cut-sheet v2): screenplay duplicate of the BSD confrontation — prose version kept in 9_7meetings.tex ============
% \section*{Eleanor \& Willard's lost focus}
% 
% \lettrine{S}{tacks} of old papers are scattered across the desk, some a little yellow, others mellow with their edges curled with age. A dim light flickers against the dusty draft Eleanor holds — an old copy of their elliptic cryptography work, marked with Willard’s equations in the margins. Eleanor stands near the desk, flipping through a brittle draft—their abandoned work on elliptic cryptography. Willard watches her, leaning against the desk, arms crossed. His eyes flick between her and the pages, an old tension creeping back into his voice.
% 
% \smallskip
% 
% Willard \textit{(narrowing his gaze, tone laced with something unspoken)} 
% 
% "I can’t believe you kept this."
% 
% \smallskip
% 
% Eleanor \textit{(not looking up, voice quieter)} 
% 
% "I didn’t. It was buried in the archives. Someone found it and left it for me."
% 
% \smallskip
% 
% Willard’s fingers perceptibly whiten where they are pressed against the desk. He exhales sharply. His \textit{(voice tightening, glancing at her)}  
% "And let me guess—you think someone’s stolen it."
% 
% \smallskip
% 
% Eleanor doesn’t answer right away. Her fingers hover over a line of proof, the one that had kept them up many-a-night all those years ago. The one they had debated, rewritten, toiled over.
% 
% \smallskip
% 
% Eleanor \textit{(softly, almost to herself)} "We were close, Willard."
% 
% \smallskip
% 
% Willard \textit{(shaking his head, voice dry)}  
% "Close doesn’t count for anything. We abandoned it, Eleanor. And now someone else is making hay with it."
% 
% \smallskip
% 
% She snaps the paper shut, finally looking at him.
% 
% \smallskip
% 
% Eleanor \textit{(tightening her grip on the pages)}  
% "We didn’t abandon it. We... lost focus."
% 
% \smallskip
% 
% Willard \textit{(a bitter laugh, arms crossing)}  
% "You lost focus."
% 
% \smallskip
% 
% The words land heavier than either of them expect. That silence of the unsaid, now thickens between them, putrid and filled with the weight of old accusations, old wounds. Eleanor holds his gaze now, making sure her voice steadier than she feels. Now \textit{(quietly, watching him)}  "You never asked me why I left."
% 
% \smallskip
% 
% Willard \textit{(bitterness dripping from him,  lowering his head)}  
% "Because I already knew. It was Olber, wasn’t it?"
% 
% \smallskip
% 
% She blinks, her breath catching for a fraction of a second. Eleanor \textit{(low, hesitant)} "You knew?"
% 
% \smallskip
% 
% A humorless chuckle escapes him. He shakes his head, staring at some point past her. Willard \textit{(with the weight of an answer he’d long stopped searching for)}  "No, but I knew it was someone. You had that look, the one people get when they want to be anywhere else. So go on, Eleanor—tell me. Was it worth it?"
% 
% \smallskip
% 
% Her fingers tighten on the paper, knuckles paling too. She holds the draft like a lifeline, like something that still might tether her to the person she used to be. Eleanor \textit{(quiet, but unshakable)}  "I was pregnant, Willard."
% 
% \smallskip
% 
% The room stills so as the industry of dust mites could be discerned. To the less discerning, muffled voices in distant corridors, the rattle and hum of a photocopier from an office down the corridor.
% 
% \smallskip
% 
% Willard \textit{(barely above a whisper, as if the words might crack if spoken too loudly)}  "What?"
% 
% \smallskip
% 
% Eleanor \textit{(finally looking up, voice steadier than her hands)}  
% "It was Olber’s."
% 
% \smallskip
% 
% The silence between them shifts, no longer heavy with accusation, but with something else — something raw, something irretrievable. Willard stares at her, searching her face for something he doesn’t quite have the words for. Maybe anger. Maybe understanding. Maybe both. Willard then \textit{(after a long pause, voice rough)}  "Why didn’t you tell me?"
% 
% \smallskip
% 
% Eleanor swallows, her grip loosening on the paper.  \textit{(quietly, almost too softly to be heard)} "I don’t know."
% 
% \smallskip
% 
% But she does. And so does he. Plans once abandoned, can never be retrieved. The weight of it lands between them. A theorem lost to time. A child that never was. An entire past neither of them ever spoke of.
% 
% \midskip
% 
% \textit{[The desk lamp hums a draft that sits between them, a relic of their their productive past, questions left unresolved.]}  
% 
%%%%%%%%%%%%%%%
% 
% ============ end commented block ============
\section*{Eleanor’s Rebuttal: Reality is Not Just a Model}

Eleanor had long tolerated the abstractionists, those who saw the Merriweather Circle as a refuge for mathematical purity, for the untainted exploration of formalism detached from the burdens of reality. She had even entertained the financial pragmatists, believing—perhaps naively — that their desire to \textit{apply} mathematics was driven by something nobler than profit. But now she sees the fractures.

\textit{They have made an error far greater than the failure of the Invariant Fund.} 

\smallskip

% \section*{Reality is Not Just an Installation of Possible Worlds}

Willard and Emily, like so many of the others, believe that mathematics dictates reality—that the world is simply one instantiation among infinite mathematically possible worlds, an accident of probability drawn from the infinite configurations of the multiverse. To them, the mathematics comes first. The world is merely one rendering.

\smallskip

\textit{But they are wrong.}\textit{Reality does not emerge from mathematical models.} Reality merely informs our mathematics. If mathematics dictated all possibilities, then every abstraction would be equally real. If our world were nothing more than an instantiation from the space of all possible worlds, then \textit{where is the law that selects this one?} Why this world, and not another? The problem is not simply one of measure, not simply of weighted probabilities or selection functions — it is a rejection of the idea that mathematics alone suffices. The world, \textit{this world}, holds \textit{clues to itself} outside of its mathematical modeling.

\smallskip

% \section*{Numerology and the Uncertainty of Discovery}

Eleanor had always argued that reality itself suggests its own mathematics—not through formal deduction, but through the recognition of patterns that cannot be derived from first principles alone. Some might call it numerology, but this was only a slur used by those who could not distinguish between coincidence and structure.

It was not pure deduction that led humans to Fibonacci spirals in sunflowers.  
It was not formalism that revealed the structure of prime numbers in biological sequences.  
It was not an axiomatic system that carved the fine-structure constant into the universe. Mathematics is not slated. \textit{It is revealed.} 

And so now, Eleanor stands at odds with those who have surrendered to pragmatism, with those who see mathematics as \textit{just} a model, rather than a process of revelation. They have wagered their future on their ability to model a system that was never meant to be modeled. She does not know yet whether this break in the Circle can be repaired. But she knows one thing with absolute certainty:

\begin{quote}
Our, one precious shared reality cannot be just a statistical anomaly in the space of all possible worlds.  
\end{quote}
\newpage
\section*{The Undercurrents of the First Light}

\lettrine{I}t was that early morning that only those gym junkies and early dippers can entertain. For all their differences, for all the fractures within the Circle, Eleanor had never thought she would see the day when Willard Jones, pragmatic, ever-calculating Willard would come back into her way of thinking. But now, as he speaks, tracing out the eigenmodes of the early universe, mapping the cosmic microwave background with spherical harmonics and oscillations of Bessel functions, she sees it: The first light of the universe carried a signature.  A signature not just of primordial sound waves, but of the very asymmetry that made them \textit{and made us}. 

\begin{figure}[h]
    \centering
    \includegraphics[width=0.8\textwidth]{Illustrations/vign-eleanorLife.png}
    \caption{Eleanor revealing her handedness to herself.}
    \label{fig:ernest_demi}
\end{figure}
% \section*{Harmonics and the Break from Perfect Order}

Willard speaks all Ernest-like of quadrupole moments, of polarization, of chiral modes that encoded structure before structure even existed. Eleanor then realizes:

\begin{quote}
These fluctuations, these first resonances of the cosmos,  are more than just polarised baryonic echoes being the break from total conformal symmetry.  They are the mathematical shadow of how we \textit{finite beings of purpose} came to be.
\end{quote}

Chirality, the handedness of the universe, was embedded in the helicity of its first light, imprinted into its fundamental vibrations. The breaking of complete order into structured imperfections allowed life, consciousness, and will to come into existence. The harmonics of the early universe were seeded with conformal asymmetry.

\begin{itemize}
    \item Polarized light in the CMB reveals a preference, a directionality in the primordial universe.
    \item Quadrupole moments encode the break from pure isotropy.
    \item Enantiomer chemistry, the handedness of molecules, ensures that life’s essential components—amino acids, sugars—exist with a \textit{fundamental chirality}.
    \item Human perception, even our grasp of mathematics, is bounded by our own physical asymmetry.
\end{itemize}

If the universe had retained perfect order, we would not exist.  
If we were simply the product of abstract mathematical models, we would not seek meaning.  

\midskip

% \section*{The Temporary Order We Cling To}

Willard is speaking in mathematics, but Eleanor hears something else. She hears something deeper, a realization that unnerves her:

\begin{quote}
We are aggregates of apparent purpose,  
fearfully aware of the finiteness of our own impermanence.
\end{quote}

She had always argued that mathematics is revealed, not imposed. But now she understands the greater truth: \textit{Mathematics does not govern reality.  Rather, reality sings mathematics into being.}

It is not we who impose structure onto the cosmos. We as the tentacles of the Big Bang have the harmonics of the cosmos \textit{imposed structure onto us}.

\smallskip

For the first time, Eleanor sees Willard differently.  Not as the pragmatist who abandoned pure theory for financial modeling. Not as the man who believed markets could be mapped onto deterministic lattices. But as someone who, whether he realizes it or not, has walked the long way around only to arrive at her argument from a different path.  

\smallskip

Upon Eleanor reality then dawns, \textit{Reality is not just a stochastic selection from possible worlds.}  
It is this world, \textit{this one} which has left imprints of its logic upon us. Our Mathematics is a consequence, not a cause. 

%%%%%%%%%%%%%%%

\section*{E \& W: a Betrayal of Minds}

It was never supposed to happen. Not like this. For all their high-minded discourse, for all the purity of their pursuit, for all their musings on purpose, mathematics, and emergent order, Eleanor and Ernest find themselves reduced to something far more primal. Desire, need, the unrelenting drive to persist as animals in a world that is indifferent to intellect.

The project—\textit{their project}—was supposed to be about revealing the deeper truths that governed reality. That intelligence could grasp, however fleetingly, the patterns beneath existence. That mathematics, symmetries, and invariants were the language of something higher. And yet, here, now, in the dark heat of their respective entanglements, they embodied the lie.

\smallskip

Eleanor had once argued that mathematics was  but a discovery, something etched into the fabric of the cosmos itself. That their work, their research, was an act of peeling back layers of illusion to reveal that truth. That there was purpose woven into the constraints that shaped their existence. But now, as she lies beside Ernest, breath unsteady, the taste of him still on her lips, she understands:  

\begin{quote}
The horror is not that there is no higher order. The horror is that they are bound to one, but it is not the one they believed in.  
\end{quote}

She had imagined herself an agent of reason, a seeker of truth, a mind elevated above the mere mechanics of biological impulse. But her body had acted before her mind could rationalize. And in that moment, she had become the very counterexample to everything she once believed.

\smallskip

Ernest, staring at the ceiling, knows it too. He has betrayed the illusion that he was ever anything other than a self-preserving agent, dressing his base impulses in the language of intellect. 

\begin{quote}
They may well have spoken of dualities, of mathematics revealing the tensions that defined existence.  And yet, here they were, ruled by the simplest, oldest inequality of them all:  
\[
\text{Desire > Ideals}
\]
\end{quote}

\smallskip

\textit{"This doesn’t change anything,"} Eleanor murmurs, though she knows it is a lie.

\textit{"Of course not,"} Ernest replies, though he knows it changes everything.

\smallskip

For all their musings on \textit{emergent intelligence}, on \textit{patterns of structure arising from chaos}, neither of them had considered the more fundamental truth:

\begin{quote}
There is no intelligence without the drive to survive.  And there is no survival without its own necessary deceptions.  
\end{quote}

Their betrayal is metaphysical being the admission that they, too, are just vehicles for persistence. That whatever noble intellectual pursuit they had framed their worth within, at bottom, the unrelenting desire to survive in some future human shape, if not in an idea, will always win. And still, they must return to the Circle. To their work. To the illusion of higher thought.

\textit{And pretend none of it ever happened.}

\newpage

\section*{The Price of Attribution}

Eleanor turns the page of an old volume, running her fingers along its softened spine, when a commotion disrupts the shop’s usual murmur of careful browsing and quiet transactions. 

A man—mid-forties, well-dressed in that unremarkable, respectable way—stands near the counter, one foot raised slightly as though it were a wounded relic. His voice, measured but forceful, carries through the narrow aisles.

\smallskip

\textit{"I stubbed my foot on that step—no warning sign. This is unacceptable."}

\smallskip

The volunteer behind the till, an elderly woman with an apologetic smile, tilts her head. \textit{"I’m terribly sorry, sir, but—"}

\smallskip

\textit{"I need to speak to the manager. A shop should take responsibility for hazards."}

\begin{figure}[h]
    \centering
    \includegraphics[width=0.8\textwidth]{Illustrations/vign-stubbMan.png}
    \caption{The third Sector is ripe for Efficiency savings.}
    \label{fig:ernest_demi}
\end{figure}

\textit{A hazard.} Eleanor glances at the step — barely a rise in the floor, an innocuous shift in level like the kind found in a hundred old buildings. No jagged edge, no real cause for injury. Yet here he is, framing his misstep as a liability. Not misfortune. Not carelessness. A claim. A potential settlement. 

\smallskip

\textit{The depths humans will plumb in order to avoid work.} 

\smallskip

The man has - or at least has had a job. A career. A comfortable enough existence to shop here not out of necessity, but as a hobbyist of forgotten books and curios. Yet at the faintest chance of \textit{extracting something for nothing}, he turns predatory, tapping a \textit{charity}—one that provides mosquito nets, medical aid, meals for those who have nothing. A sudden nausea overcomes Eleanor. It is not the injury she disbelieves — it is the immediate instinct to \textit{monetize} it. How thin the veneer of respectability must be, how close to the surface this impulse lies.

\smallskip

She wonders if he sees it in himself. She doubts it. He will frame this as \textit{holding an institution accountable}. As \textit{ensuring public safety}. He will not let himself believe what Eleanor sees so clearly: 

\smallskip

\textit{A man who would rather fight to be paid for his inconvenience than work to earn his keep.} A sharp contrast to the book she still holds in her hand. Words, worn through a dozen owners, a book purchased three, four times over, each time funding a net that might keep a child from dying of malaria.

\smallskip

\textit{What price a nameless book, a nameless idea, passed from hand to hand? What price a step misplaced?}

\smallskip

The manager appears, shaking hands, offering an appeasement. Eleanor does not stay to hear the resolution. Some things are not worth knowing. Eleanor drifts further through the narrow aisles of the charity bookshop, her fingers grazing the spines of books long divorced from their original owners. The scent of old paper, all musty, comforting, not in the least bit transient fills the air. Each book here has lived many lives, passed through many hands, yet none will ever know its journey.

\smallskip

Her idle browsing is interrupted by a voice at the counter. An older gentleman, neatly dressed but with a practiced air of frugality, is arguing with the volunteer. He gestures toward a battered book ... some trifling economic treatise, its cover worn, its margins undoubtedly filled with past readers’ feintly tragic addenda.

\smallskip

\textit{"I buy books here all the time,"} he insists. \textit{"Surely I can get a discount?"}

\smallskip

Eleanor stiffens.

\smallskip

\textit{A discount?} On a book whose proceeds will buy mosquito nets, feed the hungry, keep the cold from shivering?

\smallskip

She wants to intervene, to remind him that the words on those pages though sold third, fourth, fifth-hand with no attribution to their original writer, still retain their purpose. That this is not a transaction in the usual sense. That he is not bargaining with a merchant, but with a cause.

\smallskip

She lets the moment pass, but it gnaws at her. \textit{What price an attribution-less purchase?}

\smallskip

Her own ideas had been purchased, after all. Not by eager second-hand readers, but by financiers. Men who, in buying intellectual capital, were acquiring her as another instrument of speculation.

\smallskip

She had wrestled with this before: the unease of knowledge being packaged, repurposed, and leveraged. But now, in this bookshop, she sees its shadow more clearly.

\smallskip

\textit{Perhaps those fund investors, with their newfound wealth, will do great things.}

\smallskip

Perhaps they, in liberating themselves from financial constraints, will find the time to think, to reflect, to contribute to something greater. Perhaps the L'Hôpitalian purchase of Bernoullian ideas was always justified—\textit{if only for the greater good that trickled down from it.}

\smallskip

Or perhaps she is merely justifying her own compromises.

\smallskip

Eleanor exhales, picks up a book, and takes it to the counter. She does not ask for a discount.

\section*{A Quiet Resignation}

Eleanor steps into the dim sanctuary of her local church, drawn not by faith, but by the solace of its silence. The wooden pew beneath her is worn smooth by centuries of hands, by the weight of prayers whispered into the hush of vaulted stone.

She exhales, closing her eyes for a moment. Here, no one debates. No one justifies. No one asks whether an idea is attributed or lost in time. 

\smallskip

And then, she sees them.

\smallskip

In the far corner, barely noticed, a young child plays. Two, maybe three years old. Moving so quietly, so delicately, as if the air itself is sacred. Her mother sits beside her, silent. Not instructing. Not distracting. Just \textit{being}.

\smallskip

Eleanor watches, transfixed. \textit{How is this possible?} In an age where everyone demands to be heard, where noise is currency, where minds clamor for acknowledgment—\textit{how can two people exist together in silence and be content?} She had always believed that her work, her ideas, were meant to be carried forward. To be spoken, debated, refined by minds yet to be born. She had agonized over whether her contributions would be credited, or if others would ride upon them unacknowledged. 

\smallskip

But watching the mother and child, she understands.

\textit{I do not need to be heard. Not for now, not for eternity.} 

\smallskip

To be here, quietly among kith and kin, is enough.

\smallskip

She lets out a slow breath. \textit{Let those who build on my work take what they will. Let those who stand on my shoulders do as they wish.} 

\smallskip

We are here for so little time. \textit{Why waste it worrying about the future minds we will never know?}

\smallskip

For the first time in a long while, Eleanor feels no need to fight. No need to justify. The quiet is enough.

\newpage

\section*{Atonement Through Mathematics}

\lettrine{W}{hen} Eleanor and Willard embark on building the Merriweather Circle, they believe it is about mathematics. About creating something enduring, a space where theorems and proofs transcend fleeting politics and financial imperatives. But in reality, it is about something deeper—something neither of them dares to articulate.

Beneath their pursuit of intellectual rigor lies an unspoken truth: they are atoning for an unrecognized loss. For a moment in time when their brilliance flickered but was snuffed out by hesitation, by betrayal, by the tangled web of personal history. They are waging a quiet war—not against ignorance, but against the ghosts of what might have been.

% \subsection*{A Buried History}

Among the stacks of yellowed notes and unfinished proofs, there is a story they refuse to acknowledge. The elliptical cryptography draft that once sat on Eleanor’s desk, abandoned, lost in the fray of human frailty. A paper never submitted, an idea never tested. The betrayal—Ernest Olber’s role in distracting Eleanor at the height of their research, his involvement in her past. The child she never had, the life Willard never knew.

They speak of the Circle as a monument to pure thought. But somewhere, buried in their equations, is a missing piece of a greater story—one that neither is yet ready to unearth.

\subsection*{The Guardian of Silence}

Demi-Moorish watches them both with quiet calculation. She alone understands the full implications of dredging up the past. The truth, the full truth, would only serve to unravel the very thing they are trying to build. She sees the fault lines running beneath the foundation of the Merriweather Circle, and she knows that some things are better left buried.

Because secrets are not just personal. They are structural. They hold together the fragile alliances that make greatness possible. Demi has buried her own past, just as Eleanor and Willard have buried theirs. Some wounds never heal—but they can be repurposed, transformed into something larger than individual regret.

And so, the Merriweather Circle will rise—not as a shrine to purity, but as a testament to endurance. To the truth that remains unspoken, but never truly forgotten.

\newpage

\section*{At home Demi and Ernest}

\lettrine{T}{he} evening air is thick with silence, broken only by the occasional sound of ice settling in their untouched glasses. A low fire flickers in the hearth, casting shadows over the room. Demi stands by the mantelpiece, running her fingers absently over the frame of an old wedding photo—her and Ernest, beaming, Sylvia at her side as maid of honor. Younger, full of hope.

Demi doesn’t have to look at Ernest to know he’s watching. His body is still, his usual sharp composure slightly unguarded.

\midskip

\textit{Demi} \textit{(voice even, without turning to face him)}  
"You’ve been quiet."

\midskip

\textit{Ernest} \textit{(exhales, rubbing his temple)}  
"Long day."

\midskip

\noindent
She turns the frame slightly, adjusting it though it doesn’t need adjusting.

\midskip

\textit{Demi} \textit{(musing, more to herself than to him)}  
"It’s funny how some faces never change in your memory. I still see her like this. Before."

\midskip

Ernest doesn’t answer.

\midskip

Demi finally looks at him. His gaze is locked on the photo—more precisely, on Sylvia’s face. The way his eyes soften, just for a moment, before something unreadable shutters them again.

Demi knows grief when she sees it. She knows guilt even better.

She lets the silence stretch before speaking carefully.

\midskip

\textit{Demi} \textit{(quietly, watching him)}  
"She doesn’t hate you, you know."

\midskip

\textit{Ernest} \textit{(jaw tensing, voice low)}  
"She should."

\midskip

Demi watches him for a long moment, weighing what she already suspects against what she’s willing to know. What she’s willing to carry.

\midskip

\textit{Demi} \textit{(softer now, moving to the sofa, a quiet understanding in her tone)} 

"Maybe. But we all choose what to bury, Ernest. Don’t we?"

\midskip

For the first time that evening, he meets her eyes. There’s understanding there. A truce, perhaps. Neither of them says anything else. The fire crackles. The past remains unspoken.

\section*{The Weight of Legacy: Ernest's Confession to Chris}

A private lounge in a city that had long since lost its luster for Ernest. The aftermath of Invariant Arbitrage LLP’s collapse had left an eerie hush over everything, as if the markets themselves were holding their breath, waiting to see who would survive and who would vanish into irrelevance. The bar was dimly lit, just enough for the amber in their glasses to glow as they sat in the corner booth. Chris had the look of a man who had stared too long into the abyss. Ernest, well—he had his own abyss, though it had taken him longer to name it.

Chris sat across from him, slumped, staring into the half-empty whiskey glass in his hand. Ernest, normally composed, was anything but. The loss of his fortune had unraveled something deep within him, and for the first time in years, he spoke without calculation, without restraint.

Chris swirled his drink, watching the liquid catch the light. \textit{"So what now?"} His voice was flat. Hollow.

Ernest exhaled, running a hand through his hair. He had meant to come here and talk about the fund, about the absolute devastation of it all—his wealth, his future, gone in the blink of an algorithm gone rogue. But that wasn’t what came out.

\textit{"Do you know what really terrifies me, Chris? Not the money—I mean, yes, that’s a catastrophe in itself. But money is transient. Always has been. No, what terrifies me is what it means for legacy."}

Chris exhaled slowly, saying nothing. He let Ernest spiral, knowing there was no stopping him now.

\textit{"You, me, Eleanor, Willard— a proper Humanist soup of reason, freedom and empathy. But in denial of an afterlife we were never going to uphold the fourth pillar, were we? The old, biological, reptilian legacy. Reproduction. No. Our kind— our over reasoned kind chases production instead. We future proof ourselves through the production of immutables not mutables, that they might survive the crumbling of everything else. We don't sow seeds in flesh, we sow them in thought."}

Ernest let out a hollow laugh, shaking his head.

\textit{"And now, look at me. I have nothing. The floundering Brans-Dicke metric affine theory with torsion coupled axions that was meant to be my monument—looks more and more corned by observation. I was betting on production, and now I have lost both. I can't leave behind anything. Not to Demi. Not to anyone."}

Chris furrowed his brow. "Demi? You think she’d leave you over this?"

Ernest hesitated.

\textit{"Demi is still young enough. If she chooses, she could still take the old path. Have a child. Pass something forward, something more durable than footnotes in an astrophysical journal. And yet, I wonder—has she built up enough work to feel safe in her own name? Or is she still caught between paths, unsure if she has done enough to be remembered by thought alone?"}

Chris tilted his glass in thought. "And what about you? Have you ruled out repo-legacy entirely?"

Ernest let out a breath, gaze distant.

\textit{"Sylvia made sure of that, didn’t she? You remember. The accident. The chair. That wasn’t just an accident, Chris. That was me. My ambition, my recklessness. I was always a step ahead of myself, always pushing further. And she—she paid the price for my brilliance. I can’t justify bringing something into the world when I failed so catastrophically at protecting what was already there."}

Chris was silent for a long moment. "And yet…"

Ernest lifted his gaze, waiting.

\textit{"Yet I saw you with Jovannah, the way you look at her. Are you sure you’ve ruled it out? Or are you simply afraid of what it would mean if you hadn’t?"}

Ernest clenched his jaw, but there was no use in denying it.

\textit{"Jovannah is possibility. She is the thing I never allowed myself to imagine. She—she is life still unfolding. Unlike us, she is not burdened by proving herself immortal through work alone."}

Chris conjured a smirk, a faux flicker of amusement amidst the ruin. "So, she terrifies you too?"

Ernest laughed, but it was a tired sound. "More than losing my fortune, Chris. More than anything."

They sat there for a while in the dim light, two men who had tried to defy time in their own ways, now confronted with how little of it they truly controlled.

For all their calculations, for all their theorems and models, legacy was still the one variable they had yet to solve.

\textit{"Jovannah is fixed on Willard."}

Chris raised an eyebrow. Not what he expected. Ernest let out a dry laugh, shaking his head.

\textit{"You knew that, didn’t you?"}

Chris shrugged. \textit{"She has a way of looking at him like he’s a puzzle she wants to solve. But then again, she looks at a lot of things that way."}

\textit{"No,"} Ernest said. \textit{"Not like that. Not like a curiosity. Like he's the only question worth answering."}

Chris didn’t argue. Maybe because he had already seen it, or maybe because he simply didn’t care.

Ernest exhaled through his nose. \textit{"I knew it too. Knew it from the moment she first sat in on one of his lectures and asked him something so specific about eigenfunctions that I almost laughed. I should have laughed. I should have walked away."} He took a long sip of whiskey, letting the burn fill his throat. \textit{"But I didn't. Because, somewhere along the way, I started thinking maybe I could have something too."}

Chris tilted his head. \textit{"And Demi?"}

Ernest set his glass down, turning it slowly between his fingers. \textit{"Demi,"} he repeated. \textit{"She’s still ovulating, you know. Not for much longer, but enough to think about it."}

Chris said nothing. He was broken himself and his half measured invocations were enough for both tonight.

\textit{"I keep telling myself,"} Ernest went on, \textit{"that I never considered repo-legacy. That I was always going to leave behind a body of work instead. Some immutable truth that outlives me. That’s what Eleanor believes in, and I’ve always thought she was right."} He let out a humorless chuckle. \textit{"But it’s different, isn’t it? It’s different when you realize you have nothing to show for it. When you’ve spent your life chasing something theoretical and the only tangible thing you could leave behind is the one thing you swore you didn’t need."}

He looked at Chris then, really looked at him, and saw something flicker behind his exhausted eyes. Recognition, maybe. Or pity.

Chris took a sip of his drink. \textit{"So that’s what this is about? A legacy crisis?"}

Ernest scoffed. \textit{"That’s rich, coming from you. What was Invariant Arbitrage, if not a bid for permanence? If not a desperate attempt to make your mark?"}

Chris didn’t flinch, but his fingers tightened around his glass. \textit{"At least I tried."}

\textit{"And look where it got you."}

Silence stretched between them. Ernest closed his eyes for a moment, feeling the exhaustion settle into his bones. Then, quietly, he admitted it.

\textit{"I don't know if I have anything left to give her."}

Chris didn’t ask who he meant—Demi, or Jovannah, or even Eleanor. Maybe it didn’t matter.

\textit{"And Eleanor?"} Chris said finally, watching him. \textit{"She still thinks she’s the one pulling the strings, doesn’t she?"}

Ernest laughed, shaking his head. \textit{"She has no idea. She thinks she knows everything, and yet she doesn’t see the simplest thing of all—that she’s just as in the dark as the rest of us. That Willard’s already found something else, and that I..."} He trailed off.

Chris finished for him. \textit{"That you were never hers to begin with."}

Ernest nodded, his throat tight. \textit{"Yeah."}

Chris set his glass down with a quiet clink. \textit{"So what do we do?"}

Ernest looked at him then, at the man who had once believed he could reshape financial reality to his will, now just another casualty of entropy.

\textit{"I don't know, Chris."} He leaned back, staring at the ceiling. \textit{"But I think we’ve already lost."}

Ernest exhaled sharply, rubbing a hand over his face. His words had run dry, his mind still tangled in the threads of Jovannah and Demi, of gravity and axions, of the brute, unavoidable truth that his grand, elegant theories did not place food on the table or secure any real future. Chris had sat through it all, offering only a sip of his drink, a steady gaze that betrayed no interruptions.  

For a moment, silence settled between them, thick and unspoken.  

Then, as if by obligation, Ernest finally asked, “And you? What of you and Emily?”  

It was an effort. Chris could tell. Not feigned, but not effortless either—like an old acquaintance offering condolences long after the funeral, knowing it changes nothing, but saying it anyway.  

Chris let out a dry, humorless chuckle.  

“What of Emily.”  

Not a question. A statement, hollow and final.  

Ernest didn’t press, only nodded slightly. The barest flicker of understanding passed between them—understanding that Emily had always been a figure more carved out of ambition than affection, that whatever the two of them had built had been transactional in a way neither had ever quite admitted aloud.  

Still, Chris found himself speaking.  

“She’s… resilient.” He turned his glass in his hands, watching the way the liquid clung to the sides before slipping back down. “She’s calculating. She’ll find a way to land on her feet.”  

A pause.  

“She always does.” 

He didn’t need to say the rest. That she had been planning for disaster long before he even saw it coming. That maybe she had even seen it coming and chosen to let it play out, adapting, adjusting, ensuring her own survival.  

That, if he were being honest, she was probably the only one of them who would walk away from this whole thing relatively unscathed.  

Ernest tilted his head, something unreadable in his gaze.  

“And you?” 

Chris let out another breath, shaking his head. “Me?” He downed the rest of his drink in one slow, deliberate motion.  “I’m not so sure.”

\newpage
\section*{The Solopocists' Circle}

\begin{itemize}[leftmargin=1.3cm]

\item[Eleanor:] They’re at it again, Nina. Folding chairs arranged like a council of lost prophets. But their only rite is... hesitation.

\item[Nina:] The Solopocists?

\item[Eleanor:] Hah—Solopocists, Climatosophers, Thermodoubtists, Certaintarians. A brotherhood bound by the sacred belief that belief itself is premature clutching unconviction like it’s a virtue.

\item[Nina:] Isn’t that harsh? Some of them are just cautious thinkers. Surely skepticism has a role in science?

\item[Eleanor:] Skepticism, yes. Paralysis masquerading as rigor? That’s another game. One of them just declared: “Before I recycle, I must first resolve the ontology of carbon.”

\item[Nina:] You’re joking.

\item[Eleanor:] Not even a little. And another, stroking his epistemological beard, mused, “How do we define fire?”—while the house burns. They nod, not in agreement—God forbid—but in mutual admiration for their stalemate.

\item[Nina:] But maybe they just want to avoid mistakes.

\item[Eleanor:] Or avoid action. There’s a difference. Certaintarians won’t commit until every glacier signs a notarized affidavit. And the Carbon Quilter cites a old letter to the editor revealing consensus as “uncertain.”

\item[Nina:] And the others let that slide?

\item[Eleanor:] They nod again. Because to disagree would require a position. But they... lead to consequences. 
% Better to linger in obtuse solidarity.

\item[Nina:] That’s tragic. But at least the others—activists, intersectional theorists—they’re trying to do something, right?

\item[Eleanor:] You mean the Identity Politic Ideologues in the second room? Oh, they’re active. Dissecting every label, defending every pronoun like a battlefield. They argue in earnest.

\item[Nina:] Isn’t that a kind of solidarity?

\item[Eleanor:] Perhaps. But it’s a fractious one. The Existidealists think visibility is liberation. The Taxonomaniacs want ever-finer distinctions. And the Insistasepartists believe empathy is appropriation.

\item[Nina:] You think they’re wrong?

\item[Eleanor:] I think while they argue over lexicons, the atmosphere thickens, the ice melts, the soil erodes—and none of that gets discussed unless refracted through identity. They fight for pieces of the human experience... but forget the human race.

\item[Nina:] So what do we do?

\item[Eleanor:] We check the third room. The quiet one. The plush one. That’s where the Happy Puppeteers sit—watching, sipping carbonated smugness. No ontology. No pronouns. Just power.

\item[Nina:] Puppeteers?

\item[Eleanor:] Billionaire-backed manipulators. They count while we debate. They smile while the Solopocists stall and the Identity Circlers splinter. Because so long as the rooms are separate... the house remains theirs.

\item[Eleanor:] We weren’t supposed to beg. That was the whole point of this Circle—autonomy, intellectual clarity. And now you want us to pitch it to a venture fund like it’s an app?

\item[Nina:] I want us to survive, Eleanor. To carve a space where we aren’t crushed by administration and endless KPIs. If that means wearing a blazer and smiling at a money manager, so be it.

\item[Eleanor:] But we built this to be above the market. Pure mathematics. Pure inquiry. The Circle was supposed to be free.

\item[Nina:] And it can be. But first it needs to exist. That means lights, bandwidth, salaries. The Solopocists debate climate while the world burns. We can’t be the mathematicians who debate purity while our postdocs ghost out of academia.

\item[Eleanor:] It still feels like selling out.

\item[Nina:] It’s buying in. To the only game being played. You can sneer at the rules, but I’d rather bend them than be left outside the arena.

\item[Eleanor:] So we smile for the patrons. Say all the right euphemisms—“innovation,” “cross-disciplinary synergy,” “ROI on abstraction.”

\item[Nina:] Yes. Then we walk away with the grant. We build the Circle. We give ourselves time to think again—really think. Not in fragments between teaching overload and impact statements.

\item[Eleanor:] You play a sharper game than I expected, Nina.

\item[Nina:] Idealism costs, Eleanor. Let's make sure we can afford it.

\item[Eleanor:] Alright, strategist. If you’re going to dress our rebellion in the language of capital... I suppose you’d better write the pitch.

\item[Nina:] You mean that?

\item[Eleanor:] I mean I’ve lost too many hours fighting gatekeepers with open letters and unmet ideals. You, at least, seem capable of gaming the system without believing in it.

\item[Nina:] I’ll need a few quotes from you. Something with teeth. Gravitas still sells—even to bankers.

\item[Eleanor:] Then let’s give them a Circle that’s harder to flatten than a spreadsheet. Show them how dangerous a free mind can be—if funded.

\item[Nina:] I’ll call it: \emph{The Circle Unbound.}

\end{itemize}

\newpage
% Page 1
\begin{tcolorbox}[colback=gray!5, colframe=black, boxrule=0.3mm, sharp corners, width=\dimexpr\textwidth+1.1cm\relax, enlarge left by=0mm]

{\scriptsize\ttfamily

\vspace{1em}
\noindent\textbf{\large Mission Statement}
\vspace{0.5em}

The \textit{Circle Unbound}, CU is a mathematicians-led initiative dedicated to reclaiming time, focus, and intellectual autonomy from the administrative overload of modern academic life. Our goal is to create a minimally governed, maximally generative research collective—a space where mathematical thought can flourish free from institutional metrics, while remaining rigorously accountable to intellectual merit.

\vspace{0.8em}
\noindent\textbf{The Problem We Address}
\vspace{0.3em}

Modern research institutions—especially in pure mathematics—are burdened by bureaucratic accretion: grant-chasing, compliance paperwork, performance frameworks. These choke the creative process and disincentivize deep, long-horizon thought. Mid-career mathematicians are drowning in service loads. Promising early-career researchers are exiting academia altogether. Meanwhile, some of the most vital mathematical advances have historically emerged from protected, minimally distracted collectives.

\vspace{0.8em}
\noindent\textbf{Our Solution}
\vspace{0.3em}

CU proposes a lean, self-directed fellowship model for mathematical researchers. We seek funding to support:
{\scriptsize
\begin{itemize}
  \setlength\itemsep{2pt}
  \setlength\parsep{0pt}
  \item A rotating cohort of 12--15 researchers across all career stages.
  \item Living stipends, health coverage, and research travel.
  \item Shared space and infrastructure, either physical or distributed.
  \item Annual review by a peer-chosen external panel.
\end{itemize}
}
\vspace{0.8em}
\noindent\textbf{Why This Matters to Funders}

{\scriptsize
\begin{itemize}
  \setlength\itemsep{2pt}
  \setlength\parsep{0pt}
  \item \textit{High-leverage investment:} Mathematics underpins innovation across AI, cryptography, epidemiology, finance. Funding foundational research often yields exponential downstream returns.
  \item \textit{Narrative capital:} Backing CU positions your party as a patron of enduring human inquiry.
  \item \textit{Resilience portfolio:} Long-horizon thinking strengthens societal foresight. CU is a bet on deep cognition in a noisy age.
\end{itemize}
}
}
\end{tcolorbox}

% Page 2
\begin{tcolorbox}[colback=gray!5, colframe=black, boxrule=0.3mm, sharp corners, width=\dimexpr\linewidth+1cm\relax
, enlarge left by=0mm]
% {\footnotesize\ttfamily
{\scriptsize\ttfamily

\vspace{0.8em}
\noindent\textbf{Governance and Deliverables}

{\scriptsize
\begin{itemize}
  \setlength\itemsep{2pt}
  \setlength\parsep{0pt}
  \item An independent advisory council drawn from academia and industry.
  \item Annual public symposium and open-access research briefs.
  \item Optional thematic sprints with aligned corporate/NGO partners.
\end{itemize}
}

\textbf{Ethical Rider:} Regardless of whether funding originates from private foundations, government bodies, or institutional consortia, \textit{public engagement is non-negotiable}. In an age where mathematical innovation underpins transformative domains such as AI and quantum computing, research outputs must be ethically anchored.

\vspace{0.8em}
\noindent\textbf{Why Government Funding is Vital}
\vspace{0.3em}

As large-scale technological platforms increasingly consolidate knowledge production, democratic governments must not outsource mathematical sovereignty to corporate monopolies. CU provides an independent reservoir of critical thinking. It lightens the research burden from debt-laden public institutions and returns high-impact intellectual capital to the commons. CU is not a retreat from rigor. It is a sanctuary for it. In an age of reactive metrics and monetized attention, we offer a test case in trust-driven, outcome-rich inquiry. We welcome a conversation with funders who see mathematics not as a cost center, but as civilization's compound interest.
% \vspace{0.8em}

\noindent\textbf{Budget Summary (3 Years)}

{\scriptsize
\begin{itemize}
  \setlength\itemsep{2pt}
  \setlength\parsep{0pt}
  \item Researcher stipends (15 fellows × 3 years): £2.7M
  \item Operations + admin (audit, review): £450K
  \item Infrastructure (workspace, IT, retreat): £300K
  \item Public engagement and reporting: £150K
  \item \textbf{Total Requested: £3.6M}
\end{itemize}
}

\centering
\textit{"Between rigor and refusal,\\
there lies a margin of freedom.\\
It begins as a question.\\
And sometimes,\\
a Circle."}

\vspace{1.5em}
\textbf{\LARGE $\bigcirc$}

\vspace{2em}
\begin{flushright}
\texttt{Respectfully submitted,}\\
\texttt{Eleanor and Nina}
\end{flushright}
}
\end{tcolorbox}

\newpage

\begin{itemize}[leftmargin=1.5cm]

\item[Eleanor:] \textit{(reading aloud from the printout)} “Mathematics not as a cost center, but as civilization’s compound interest…” Hah. That’s quite a closer.

\item[] \textit{(She lowers the pages, her eyes narrowing in thought.)}

\item[Eleanor:] You sharp little tactician. You actually did it. Turned an argument for epistemic refuge into a funding pitch that sings like a policy manifesto. With deliverables.

\item[] \textit{(She scans again, lips moving silently through ‘trust-driven, outcome-rich inquiry.’)}

\item[Eleanor:] You’ve mapped the Circle onto the terrain of statecraft without losing its soul. I wouldn’t have managed that, not without barbed edges or footnotes.

\item[] \textit{(Pausing, she presses a palm to the table—half in admiration, half in reluctant surprise.)}

\item[Eleanor:] Nina, you might be the future I was too tired to expect. And just idealistic enough to make cynicism inefficient.

\item[] \textit{(She folds the document carefully, as if it were a first edition.)}

% \item[Eleanor:] Let’s take this to the people who still sign the cheques. And remind them: the mathematicians have not yet gone quiet.

\end{itemize}

\newpage
\section*{Westminster Scientific Research Subcommittee, Room 4B}

\begin{figure}[h]
    \centering
\includegraphics[width=0.8\textwidth]{Illustrations/vign-subcomittee.png}
    \caption{Scientific committee coming to reason.}
    \label{fig:ernest_demi}
\end{figure}

\begin{itemize}[leftmargin=3.5cm]

\item[Chair (Baroness Kingsley):] Item 7—submission from The Circle Unbound. A... mathematical fellowship, is it? Sounds a bit utopian. Who’s briefing?

\item[Aide:] \textit{(murmuring)} It's a non-institutional model. Small, rigorous, independent. Autonomy for mathematicians. No bureaucracy.

\item[Geoffrey Molesworth MP:] \textit{(glancing at the folder)} They’re asking for what—three and a half million? To do... less?


\item[Lord Finley of Merioneth:] No, Geoff. To think. Quietly. Persistently. There’s something... bracing in this. “Mathematics as compound interest for civilisation.” Rather fine.

\item[Dr Bhatt (STEM Advisor):] It’s surprisingly current. Their rider on public engagement and IP transparency—very prescient. Especially with the current AI licensing fracas. They're not naïve.

\item[Baroness Kingsley:] \textit{(nodding slowly)} It's not cloaked in technocratic promises. No dashboard of deliverables. Just intellectual resilience.

\item[Molesworth:] Still feels romantic. But I can’t help thinking... we fund start-ups on thinner narratives.

\item[Lord Finley:] And they request less than a single procurement error. It's lean. But its reach might exceed its spend.

\item[Saira Bhatt:] They’re right about something else too. If we don’t fund the thinkers, someone else will—and it won’t be us reaping the dividends of foundational maths.

\item[Kingsley:] \textit{(closing the folder)} I propose a pilot grant. Limited term. Peer-reviewed. But with breathing room. Let’s see what minds do when untethered.

\item[Molesworth:] \textit{(grinning wryly)} Well then—looks like we’ll be investing in thought. Fancy that.

\item[] \textit{(Murmurs of assent. A rare moment of accord. And somewhere in the deep stacks of Committee Room 4B, hope makes a quiet scribble in the margin.)}

\end{itemize}

\newpage
\subsection*{Finale?: Eleanor’s Kitchen, Early Morning Light}

\begin{itemize}[leftmargin=1.5cm]

\item[] \textit{Steam curls from mismatched mugs. The kettle ticks as it cools. Papers are strewn across the kitchen table—handwritten notes, grant drafts, old theorems revisited.}

\item[Eleanor:] \textit{(reading the email aloud)} “The Committee has agreed to fund an initial pilot. Terms modest, but independent. Disbursement pending.” Hmph. They blinked.

\item[Nina:] \textit{(leaning against the counter)} No ceremony. No ribbon. But they blinked.

\item[Eleanor:] It’s not enough. You know that.

\item[Nina:] It’s never enough. But it’s time. Space. Sanctioned subversion.

\item[] \textit{(Eleanor folds the letter back into its envelope—carefully, as if handling something more delicate than it looks.)}

\item[Eleanor:] It’s a splinter. In a world of vast engines and soft tyrannies, we’ve won... a splinter.

\item[Nina:] But it's ours. And if we wedge it just right—

\item[Eleanor:] \textit{(smiling despite herself)} The door might creak open?

\item[Nina:] Or at least make a satisfying noise.

\item[] \textit{(They both look out the window, where the street is still grey, the light muted. But in the silence, a certain clarity hums—a note that lingers longer than expected.)}

\item[Eleanor:] Well then. Let’s begin.

\end{itemize}







